\documentclass[12pt]{article}
\usepackage{amsmath, amssymb}
\usepackage[margin=1in]{geometry}
\usepackage{parskip}

\title{ORF 309 -- Homework 4, Q1}
\author{}
\date{}

\begin{document}
\maketitle

\section*{Q1(a): Compute $P\{X + Y = k\}$}

\textbf{Setup.} Let $X \sim \text{Binom}(n, p)$ and $Y \sim \text{Binom}(m, p)$ be independent random variables.

\textbf{Probabilistic argument.} $X$ counts the number of successes in $n$ independent trials, each with success probability $p$. Similarly, $Y$ counts successes in $m$ independent trials with the same probability $p$. Since $X$ and $Y$ are independent, $X + Y$ counts the total number of successes across all $n + m$ independent trials. Therefore:
\[
    X + Y \sim \text{Binom}(n + m,\, p)
\]

\textbf{Result.}
\[
    \boxed{P\{X + Y = k\} = \binom{n+m}{k}\, p^k\, (1-p)^{n+m-k}, \quad k = 0, 1, \ldots, n+m}
\]

\textbf{Verification via convolution.} One can also verify this mechanically using the law of total probability. Since $X$ and $Y$ are independent:
\begin{align*}
    P\{X + Y = k\} &= \sum_{j=0}^{k} P\{X = j\}\, P\{Y = k - j\} \\[6pt]
    &= \sum_{j=0}^{k} \binom{n}{j} p^j (1-p)^{n-j} \cdot \binom{m}{k-j} p^{k-j} (1-p)^{m-(k-j)} \\[6pt]
    &= p^k (1-p)^{n+m-k} \sum_{j=0}^{k} \binom{n}{j}\binom{m}{k-j}
\end{align*}

By the Vandermonde identity, $\displaystyle\sum_{j=0}^{k} \binom{n}{j}\binom{m}{k-j} = \binom{n+m}{k}$, which recovers the same answer.

\section*{Q1(b): Compute $P\{X = Z - 1\}$}

\textbf{Setup.} Let $X \sim \text{Binom}(n, p)$ and $Z \sim \text{Geom}(p)$ be independent. Recall that $Z$ is the number of trials until the first success, so $Z - 1$ is the number of failures before the first success:
\[
    P\{Z - 1 = k\} = (1-p)^k \cdot p, \quad k = 0, 1, 2, \ldots
\]

\textbf{Calculation.} By independence, we sum over all values $k$ that $X$ and $Z-1$ can share (i.e., $k = 0, 1, \ldots, n$):
\begin{align*}
    P\{X = Z - 1\} &= \sum_{k=0}^{n} P\{X = k\} \cdot P\{Z - 1 = k\} \\[6pt]
    &= \sum_{k=0}^{n} \binom{n}{k} p^k (1-p)^{n-k} \cdot (1-p)^k \cdot p \\[6pt]
    &= p\,(1-p)^n \sum_{k=0}^{n} \binom{n}{k} p^k
\end{align*}

By the binomial theorem, $\displaystyle\sum_{k=0}^{n} \binom{n}{k} p^k \cdot 1^{n-k} = (1 + p)^n$. Therefore:
\begin{align*}
    P\{X = Z - 1\} &= p\,(1-p)^n\,(1+p)^n \\[6pt]
    &= p\,\bigl[(1-p)(1+p)\bigr]^n
\end{align*}

\textbf{Result.}
\[
    \boxed{P\{X = Z - 1\} = p\,(1 - p^2)^n}
\]

\section*{Appendix: The Binomial Theorem}

\textbf{Statement.} For any $a$, $b$, and positive integer $n$:
\[
    (a + b)^n = \sum_{k=0}^{n} \binom{n}{k} a^k \, b^{n-k}
\]

\subsection*{Proof (combinatorial)}

$(a+b)^n$ means multiplying $n$ copies of $(a+b)$ together:
\[
    \underbrace{(a+b)(a+b)(a+b)\cdots(a+b)}_{n \text{ factors}}
\]
When you expand this product, each term comes from choosing either $a$ or $b$ from each of the $n$ factors. A term with $a^k b^{n-k}$ arises when you pick $a$ from exactly $k$ factors and $b$ from the remaining $n-k$ factors.

How many ways can you pick $a$ from exactly $k$ out of $n$ factors? That is $\binom{n}{k}$ (``$n$ choose $k$'').

So the coefficient of $a^k b^{n-k}$ is $\binom{n}{k}$, and summing over all possible values of $k$ gives:
\[
    (a+b)^n = \binom{n}{0}a^0 b^n + \binom{n}{1}a^1 b^{n-1} + \cdots + \binom{n}{n}a^n b^0 = \sum_{k=0}^{n}\binom{n}{k}a^k b^{n-k} \qquad \square
\]

\subsection*{Concrete numerical example}

Let's verify with $a = 1$, $b = 2$, $n = 3$. We know $(1+2)^3 = 3^3 = 27$.

Now expand using the theorem:
\begin{align*}
    (1+2)^3 &= \binom{3}{0}1^0 \cdot 2^3 + \binom{3}{1}1^1 \cdot 2^2 + \binom{3}{2}1^2 \cdot 2^1 + \binom{3}{3}1^3 \cdot 2^0 \\[6pt]
    &= 1 \cdot 8 \;+\; 3 \cdot 4 \;+\; 3 \cdot 2 \;+\; 1 \cdot 1 \\[6pt]
    &= 8 + 12 + 6 + 1 = 27 \quad \checkmark
\end{align*}

\subsection*{How it was used in Q1(b)}

We had the sum:
\[
    \sum_{k=0}^{n} \binom{n}{k} p^k
\]
This matches the binomial theorem with $a = p$ and $b = 1$:
\[
    \sum_{k=0}^{n} \binom{n}{k} p^k \cdot 1^{n-k} = (p + 1)^n = (1+p)^n
\]
This is the step that collapsed the sum into the closed-form expression.

\end{document}
